\documentclass[12pt,a4paper]{article}
\usepackage[utf8]{vietnam}
\usepackage[T5]{fontenc}
\usepackage[margin=2cm]{geometry}
\usepackage{amsmath,amssymb}
\usepackage{booktabs}
\usepackage{tabularx}
\usepackage{array}
\usepackage{enumitem}
\usepackage{listings}
\usepackage{xcolor}
\usepackage{hyperref}
\hypersetup{colorlinks=true,linkcolor=blue,urlcolor=blue}
\lstset{basicstyle=\ttfamily\small,columns=fullflexible,frame=single,breaklines=true}

\title{\textbf{BÀI GIẢI BÀI TẬP 3\\ THUẬT GIẢI A*}}
\author{}
\date{}

\begin{document}
\maketitle
\tableofcontents
\newpage

\section{Bài 1: Tìm đường từ Arad đến Bucharest}

\subsection{Cơ sở của thuật toán}

Với mỗi thành phố $n$, thuật toán A* sử dụng hàm đánh giá
\[
    f(n)=g(n)+h(n),
\]
trong đó $g(n)$ là chi phí thực tế từ Arad đến $n$, còn $h(n)$ là khoảng
cách đường chim bay từ $n$ đến Bucharest. Ở mỗi vòng lặp, ta lấy khỏi
OPEN đỉnh có $f$ nhỏ nhất. Nếu có hòa, lời giải này ưu tiên $g$ nhỏ hơn;
nếu vẫn hòa thì ưu tiên đỉnh được đưa vào OPEN trước. Ta dừng khi
Bucharest được lấy ra khỏi OPEN.

Heuristic đường chim bay không vượt quá chi phí ngắn nhất còn lại đến
Bucharest (admissible), vì vậy A* trả về đường đi tối ưu khi dừng theo
quy tắc trên.

\subsection{Các bước thực hiện}

Các cạnh cần dùng trong quá trình tìm kiếm là
\[
\begin{aligned}
&Arad\mathbin{-}Sibiu=140,\quad Sibiu\mathbin{-}Rimnicu\ Vilcea=80,\\
&Rimnicu\ Vilcea\mathbin{-}Pitesti=97,\quad Pitesti\mathbin{-}Bucharest=100,\\
&Sibiu\mathbin{-}Fagaras=99,\quad Fagaras\mathbin{-}Bucharest=211,\\
&Rimnicu\ Vilcea\mathbin{-}Craiova=146,\quad Pitesti\mathbin{-}Craiova=124.
\end{aligned}
\]

Bảng diễn biến chính của A*:

\begin{center}
\renewcommand{\arraystretch}{1.25}
\begin{tabular}{c l r r r p{6.2cm}}
\toprule
Bước & Đỉnh lấy ra & $g$ & $h$ & $f$ & Các cập nhật quan trọng\\
\midrule
0 & Arad & 0 & 366 & 366 & Zerind $(75,374,449)$; Timisoara $(118,329,447)$; Sibiu $(140,253,393)$.\\
1 & Sibiu & 140 & 253 & 393 & Rimnicu Vilcea $(220,193,413)$; Fagaras $(239,178,417)$; Oradea $(291,380,671)$.\\
2 & Rimnicu Vilcea & 220 & 193 & 413 & Pitesti $(317,98,415)$; Craiova $(366,160,526)$.\\
3 & Pitesti & 317 & 98 & 415 & Bucharest $(417,0,417)$. Đường qua Pitesti đến Craiova có $g=441$, không tốt hơn $366$.\\
4 & Fagaras & 239 & 178 & 417 & Bucharest đã có trong OPEN với $g=417$; đường mới qua Fagaras có $g=450$, không tốt hơn.\\
5 & Bucharest & 417 & 0 & 417 & Đạt mục tiêu; dừng và truy vết con trỏ cha.\\
\bottomrule
\end{tabular}
\end{center}

Trong bảng, mỗi bộ ba có dạng $(g,h,f)$. Chẳng hạn, sau khi mở rộng
Arad, Sibiu được chọn vì $f(Sibiu)=140+253=393$ nhỏ nhất. Sau đó:
\[
\begin{aligned}
 f(Rimnicu\ Vilcea)&=220+193=413,\\
 f(Pitesti)&=317+98=415,\\
 f(Bucharest)&=417+0=417.
\end{aligned}
\]

\subsection{Kết quả}

Lần theo con trỏ cha từ Bucharest về Arad, ta được
\[
\boxed{Arad\;\longrightarrow\;Sibiu\;\longrightarrow\;Rimnicu\ Vilcea
\;\longrightarrow\;Pitesti\;\longrightarrow\;Bucharest.}
\]

Tổng độ dài đường đi là
\[
    140+80+97+100=\boxed{417\text{ km}}.
\]
Đường trực tiếp qua Fagaras có độ dài $140+99+211=450$ km, lớn hơn
417 km, nên không phải nghiệm tối ưu.

\section{Bài 2: Giải 8-puzzle bằng A*}

\subsection{Mô hình hóa}

Một trạng thái được biểu diễn bởi tuple gồm 9 số theo thứ tự từng hàng;
số $0$ biểu diễn ô trống. Mỗi nút tìm kiếm lưu trạng thái, chi phí $g$,
heuristic $h$, nút cha và nước đi đã sinh ra nó. Mỗi nước đi có chi phí
1, nên chi phí của lời giải chính là số nước đi.

Chương trình trong tệp \texttt{astar\_8puzzle.py} đọc tệp \texttt{INPUT.txt}
(theo mẫu của đề), hoặc nhận tên tệp qua dòng lệnh. Tệp đầu vào gồm một
dòng kích thước 3, ba dòng trạng thái đầu và ba dòng trạng thái đích.

\subsection{Heuristic Manhattan}

Với mỗi quân $t\ne0$, gọi $(r_t,c_t)$ là vị trí hiện tại và
$(r_t^*,c_t^*)$ là vị trí của quân đó trong mục tiêu. Heuristic được dùng
là
\[
 h(s)=\sum_{t\ne0}\left(|r_t-r_t^*|+|c_t-c_t^*|\right).
\]
Mỗi nước đi chỉ làm giảm hoặc tăng tổng khoảng cách một lượng không quá
1, do đó Manhattan là heuristic admissible và consistent. A* graph search
vì thế tìm được lời giải có số bước nhỏ nhất.

\subsection{Kiểm tra dữ liệu và tính khả giải}

Chương trình kiểm tra kích thước bằng 3, đủ sáu dòng trạng thái, mỗi trạng
thái chứa đúng một hoán vị của các số $0,1,\ldots,8$. Với bàn cờ $3\times3$,
số nghịch thế được tính sau khi bỏ số 0. Trạng thái đầu và trạng thái đích
khả biến đổi qua lại khi và chỉ khi hai số nghịch thế có cùng tính chẵn lẻ.
Nếu khác parity, chương trình in thông báo không có lời giải và không chạy
A*.

\subsection{Các bước của A*}

\begin{enumerate}[label=\arabic*.]
    \item Đưa trạng thái đầu vào priority queue với $g=0$ và $f=g+h$.
    \item Lấy trạng thái có $f$ nhỏ nhất; bỏ qua bản ghi cũ nếu nó không còn
    là chi phí tốt nhất đã biết cho trạng thái đó.
    \item Nếu trạng thái là mục tiêu, truy vết các con trỏ cha.
    \item Tìm vị trí ô 0 và thử các hướng Lên, Xuống, Trái, Phải nếu hợp lệ.
    \item Với trạng thái mới có $g$ nhỏ hơn giá trị đã lưu, cập nhật
    $best\_g$, cha và đưa trạng thái vào OPEN.
    \item Đảo ngược chuỗi truy vết rồi xuất từng ma trận trung gian.
\end{enumerate}

Quy ước mã nước đi là $U$ (Lên), $D$ (Xuống), $L$ (Trái), $R$ (Phải),
theo hướng di chuyển của ô trống. Tệp \texttt{OUTPUT.txt} chứa trạng thái
đầu, trạng thái đích, từng bước cùng các giá trị $g,h,f$, số bước và số
nút đã mở rộng.

\subsection{Mã nguồn cốt lõi}

Mã nguồn đầy đủ được bàn giao trong tệp \texttt{astar\_8puzzle.py}. Các
hàm chính gồm:
\begin{itemize}
    \item \texttt{read\_input}, \texttt{parse\_state}: đọc và kiểm tra input;
    \item \texttt{inversion\_count}, \texttt{is\_solvable}: kiểm tra parity;
    \item \texttt{manhattan}: tính heuristic;
    \item \texttt{neighbors}: sinh trạng thái kế tiếp;
    \item \texttt{a\_star}: tìm kiếm bằng heap tối thiểu;
    \item \texttt{reconstruct\_path}, \texttt{format\_solution}: truy vết và xuất kết quả.
\end{itemize}

\subsection{Độ đúng đắn và độ phức tạp}

Manhattan không vượt quá chi phí thật và nhất quán; vì mỗi cạnh có chi phí
1, A* luôn lấy nghiệm có số bước nhỏ nhất khi trạng thái đích được lấy ra
khỏi OPEN. Với 8-puzzle chỉ có tối đa $9!$ hoán vị, trong thực tế số trạng
thái được mở rộng nhỏ hơn nhiều nhờ heuristic và bảng $best\_g$. Độ phức tạp
phụ thuộc số trạng thái đã sinh, thường viết là $O(b^d)$ trong trường hợp
xấu nhất, với $b\le4$ và $d$ là độ sâu nghiệm.

\section{Cách chạy và kiểm thử}

Đặt các tệp \texttt{astar\_8puzzle.py} và \texttt{INPUT.txt} cùng thư mục,
sau đó chạy:
\begin{lstlisting}
python astar_8puzzle.py
\end{lstlisting}
Có thể chỉ định tệp vào/ra:
\begin{lstlisting}
python astar_8puzzle.py INPUT.txt OUTPUT.txt
\end{lstlisting}

Nên kiểm thử thêm ba trường hợp: trạng thái đầu trùng trạng thái đích (0
bước), trạng thái đổi chỗ hai quân (không khả giải), và input sai định dạng.

\end{document}
